\documentclass[12pt]{article}
\usepackage[T1]{fontenc}
\usepackage[utf8]{inputenc}
\usepackage{lmodern}
\usepackage{textcomp}
\usepackage{enumitem}
\usepackage{geometry}
\usepackage{graphicx}
\usepackage{amsmath, amssymb}
\geometry{margin=1in}
\begin{document}
\begin{center}
Philosophy - Undergraduate Level Assessment 14 \
Total Marks: 40 \
Answer all questions.
\end{center}

\section*{Multiple Choice (10 marks)}
\begin{enumerate}[label=\arabic*.]
\item According to the Ideal Moral Code theory, one is obligated to do what the ideal moral rules would require:
\begin{enumerate}[label=(\alph*)]
    \item in a world where moral rules are constantly changing.
    \item in a world in which everyone complied with those rules perfectly.
    \item in an ideal institutional setting.
    \item in a world where every individual creates their own rules.
    \item in a setting where institutions are corrupted.
\end{enumerate}
\item Kamm claims that Sandel's moral distinction between treatment and enhancement assumes that
\begin{enumerate}[label=(\alph*)]
    \item enhancement is always morally superior to treatment.
    \item enhancement is inherently unnatural.
    \item human beings have a moral duty to enhance themselves.
    \item medical treatment is always more morally acceptable than enhancement.
    \item there is no moral difference between treatment and enhancement.
\end{enumerate}
\item By "animal motion," Hobbes means:
\begin{enumerate}[label=(\alph*)]
    \item instinctive behavior, such as nursing young.
    \item irrational behavior.
    \item involuntary operations such as heartbeat and breathing.
    \item the physical actions and reactions of animals.
    \item human behavior that is driven by basic needs.
\end{enumerate}
\item On McGregor's view, our rights carve out
\begin{enumerate}[label=(\alph*)]
    \item the boundaries of our moral obligations.
    \item the domain of our personal identity.
    \item all of the above.
    \item what we are free to do.
    \item the responsibilities we have towards others.
\end{enumerate}
\item According to Parfit, the obligation to give priority to the welfare of one's children is:
\begin{enumerate}[label=(\alph*)]
    \item agent-relative.
    \item agent-neutral.
    \item absolute.
    \item none of the above.
\end{enumerate}
\end{enumerate}

\section*{True / False (10 marks)}
\textit{Answer True or False}
\begin{enumerate}[label=\arabic*.]
\setcounter{enumi}{5}
\item True or False: According to Kant, self-interest is the only thing that can be considered good without qualification.
\item True or False: Gauthier argues that emotional justification is essential for practical rationality.
\item True or False: Reiman asserts that van den Haag's arguments in favor of the death penalty follow common sense reasoning.
\item True or False: Reiman argues that van den Haag's position leads to the conclusion that we should institute death by torture.
\item True or False: Moral nihilism is the view that moral considerations do not apply to war, asserting that moral values are ultimately unfounded.
\end{enumerate}

\section*{Long Form Response (20 marks)}
\textit{Answer each question with a clear and complete explanation.}
\begin{enumerate}[label=\arabic*.]
\setcounter{enumi}{10}
\item Discuss Thomson's revised violinist case in which the violinist requires your kidneys for one hour. Analyze her stance on the implications for the right to use another's body.
\item Discuss the significance of the construction of the first Buddhist temple in Japan in 596 CE and its impact on Japanese philosophy and culture.
\end{enumerate}

\vfill
\noindent\textit{End of Paper}
\end{document}